\section{Въведение}
\label{sec:introduction}

\subsection{Постановка на проблема}

Достъпът до данни остава едно от местата, в които разминаването между езика на приложението и модела на съхранение създава най-много структурни грешки. В традиционните слоеве за достъп до бази данни заявките се описват като низове, които се валидират едва при изпълнение. Това означава, че неправилни имена на колони, несъвместими типове, невалидни \code{JOIN} конструкции или неправилно подадени параметри се откриват късно --- често върху реална база данни и след допълнителни ръчни тестове. Същевременно смяната на свързващия компонент от релационен към документен, графов или ключ-стойност модел обикновено изисква пълно пренаписване на слоя за достъп до данни.

Обектно-релационното картографиране е класическият подход за свързване на обектния модел на приложението с персистентното съхранение~\cite{fowler2002}. Повечето реализации на обектно-релационно съпоставяне (ОРС) обаче остават силно зависими от механизми по време на изпълнение: SQL низове, кодогенерация, анотации, макроси или динамични описания на схема. В контекста на C++ това създава допълнителен проблем. Езикът предлага мощна система от шаблони, \code{constexpr} изчисления и \code{concepts}, но тези механизми често не се използват докрай в библиотеките за достъп до данни. Настоящата работа изследва доколко е възможно компилаторът да поеме ролята на първи валидатор на структурата на заявката и на съвместимостта със свързващия компонент.

Съществена особеност на разработваната система е, че въпреки името библиотека за обектно-релационно съпоставяне по време на компилация, архитектурата ѝ не е ограничена само до релационния модел. Напротив, тя използва еднакъв C++ модел на данните и еднакво междинно представяне на заявките като логически слой над релационни и нерелационни системи. Това поставя втори изследователски въпрос: кои части на ОРС абстракцията са универсални и кои неизбежно трябва да се ограничават чрез базиран на възможности договор на свързващия компонент.

Третият проблемен пласт е свързан с начина на изпълнение. Дори когато логическият модел и междинното представяне на заявките са типово коректни, практическата стойност на една съвременна библиотека за достъп до данни зависи и от това дали тя може да интегрира асинхронно изпълнение без разпадане на типовите гаранции и без принудително преминаване към стил с функции за обратно извикване (функция за обратно извикване). Така дипломната работа се организира около две допълващи се оси: моделиране по време на компилация и асинхронна архитектура на изпълнението.

\subsection{Цел и задачи на разработката}

Основната цел на дипломната работа е да представи проектирането и реализацията на библиотека за обектно-релационно съпоставяне по време на компилация за C++, която:

\begin{enumerate}[label=\arabic*.]
    \item моделира данните и заявките чрез статично типизирани C++ конструкции;
    \item позволява изграждане на \code{SELECT}, \code{INSERT}, \code{UPDATE} и \code{DELETE} заявки като \code{constexpr} обекти;
    \item отделя логическата структура на заявката от специфичното ѝ материализиране за дадения свързващ компонент;
    \item поддържа релационни и нерелационни конектори чрез \code{connector\_trait<DB>};
    \item използва механизъм за ограничаване на неподдържани операции по време на компилация въз основа на възможностите на компонента;
    \item демонстрира как C++ моделът определя логическата схема и физическата структура на данните при различни свързващи компоненти;
    \item реализира базиран на корутини асинхронен слой, който надгражда синхронния модел на конектора по формален и типово безопасен начин;
    \item формализира договор за добавяне на нови конектори, включително когато те поддържат платформено-специфично асинхронно изпълнение;
    \item валидира архитектурата чрез систематични модулни и интеграционни тестове и чрез емпирична сравнителна (сравнителен тест) оценка.
\end{enumerate}

За постигане на тази цел работата решава следните задачи: анализ на съществуващи решения; дефиниране на статичен модел на същностите; изграждане на междинно представяне на заявкитете; проектиране на свързващ слой и labelи за възможности; изграждане на базиран на корутини асинхронен слой; имплементиране на конектори за различни типове бази данни и различни модели на изпълнение; изследване на миграции, подготвени заявки, нишкова безопасност и асинхронно управление на връзки; формулиране на критерии за разширяемост; изграждане на методология за сравнителен анализ за оценка на синхронния и асинхронния режим.

\subsection{Обхват и ограничения}

Работата разглежда библиотеката като слой за моделиране, изграждане и изпълнение на заявки. В обхвата ѝ попадат моделът на същностите, междинното представяне на заявките, механизмът за заместители (заместители), подготвените заявки, миграционният модел, архитектурата на свързващия компонент, ограничаването въз основа на възможности, базираният на корутини асинхронен слой и конкретните реализации на свързващи компоненти в хранилището. Съществена част от оценяването е сценарийният анализ на SQLite, PostgreSQL, MySQL, MongoDB, Redis, Neo4j и Cassandra, както и сравнителната оценка на синхронния и асинхронния режим на изпълнение.

Извън непосредствения обхват остават оптимизации на производствено ниво като разпределено управление на връзки, автоматизирани стратегии за миграция за всички свързващи компоненти, измерване на производителността в реално време върху пълна матрица от производствени среди и пълна експлоатация на бъдещата функционалност за отражение (отражение) на C++26. Тези направления са разгледани като перспектива в Глава~\ref{sec:future_work}.

\subsection{Структура на изложението}

Дипломната работа е организирана по следния начин.

\textbf{Глава~\ref{sec:existing_solutions}} разглежда съществуващи C++ решения за достъп до бази данни и подходи за обектно-релационно съпоставяне и позиционира разработката спрямо тях.

\textbf{Глава~\ref{sec:theoretical_background}} представя теоретичните основи на обектно-релационното съпоставяне по време на компилация и моделите за съхранение, необходими за по-късния анализ на релационни и нерелационни сценарии.

\textbf{Глава~\ref{sec:coroutines_async_theory}} въвежда теоретичния апарат на корутините и асинхронното изпълнение в C++, който е необходим за разбирането на втория основен принос на библиотеката.

\textbf{Глава~\ref{sec:architecture}} описва архитектурата на библиотеката за обектно-релационно съпоставяне по време на компилация: описанията на същностите, междинното представяне на заявките, \code{connector\_trait}, механизмите за възможности и мястото на миграциите в общия модел.

\textbf{Глава~\ref{sec:async_execution_architecture}} представя асинхронната архитектура на изпълнението и показва как базираният на корутини слой се свързва със синхронния програмен интерфейс на конектора, с модела на набор от нишки и с управлението на връзките.

\textbf{Глава~\ref{sec:backend_scenarios}} анализира как една и съща логическа абстракция и различни стратегии за изпълнение се материализират в релационни, документни, ключ-стойност, графови и широки колонни свързващи компоненти.

\textbf{Глава~\ref{sec:verification_evaluation}} обединява верификационните доказателства и емпиричната оценка чрез резултатите от сравнителните тестове.

\textbf{Глава~\ref{sec:verification_extensibility}} формализира договора за надграждане чрез нови конектори, включително изискванията за възможности и асинхронност.

\textbf{Глава~\ref{sec:future_work}} очертава ограниченията на текущата версия и възможните направления за бъдещо развитие, а \textbf{Глава~\ref{sec:conclusion}} обобщава резултатите и приносите.
